\section{Bott--Chern and Aeppli cohomology}
\label{sec:bott-chern}

A closed form of pure type can be exact for $d$ without being
exact for $\partial\bar\partial$.
The two differentials retain this distinction even when all
intermediate de Rham groups vanish. We determine their additive
cohomologies, construct the two classes in bidegree $(1,1)$,
and use the map to the base curve to compute all their products.
Integration then determines the Aeppli module.

Throughout this section, $Y$ is the compact complex threefold of
Theorem~\ref{cmp:thm:comparison}. We retain its fixed nonzero
contraction parameter, integral lifts, logarithmic increments
and adapted gluing. The Hodge calculation
\eqref{hr:eq:hodge-table} and the Betti comparison
\eqref{hr:eq:betti} concern this same $Y$.

\subsection{Two quotients of the form complex}

Write $A^{p,q}=\Gamma(Y,\mathcal A_Y^{p,q})$ for the smooth
complex $(p,q)$-forms, with $A^{p,q}=0$ outside
$0\le p,q\le3$. The differentials have bidegrees $(1,0)$ and
$(0,1)$ and satisfy
\[
 \partial^2=\bar\partial^2=0,\qquad
 \partial\bar\partial+\bar\partial\partial=0.
\]
The \emph{Bott--Chern} and \emph{Aeppli} groups are
\begin{align}
 H_{BC}^{p,q}(Y)
 &=\frac{\ker(\partial:A^{p,q}\to A^{p+1,q})
                 \cap\ker(\bar\partial:A^{p,q}\to A^{p,q+1})}
          {\im(\partial\bar\partial:A^{p-1,q-1}\to A^{p,q})},
       \label{bc:eq:definition-BC}\\
 H_A^{p,q}(Y)
 &=\frac{\ker(\partial\bar\partial:A^{p,q}\to A^{p+1,q+1})}
          {\partial A^{p-1,q}+\bar\partial A^{p,q-1}}.
       \label{bc:eq:definition-A}
\end{align}
Anticommutation puts each denominator in its numerator.
Let $h_{BC}^{p,q}$ and $h_A^{p,q}$ denote their complex dimensions.
We also use the row and column cohomologies
\[
 H_\partial^{p,q}=H^p(A^{\bullet,q},\partial),\qquad
 H_{\bar\partial}^{p,q}=H^q(A^{p,\bullet},\bar\partial).
\]
The latter is the ordinary Dolbeault group
$H^q(Y,\Omega_Y^p)$.

There are two spectral sequences. The filtrations by the first
and second bidegrees on $(\operatorname{Tot}A,d)$, with
$d=\partial+\bar\partial$, give
\begin{equation}\label{bc:eq:two-sequences}
 \begin{aligned}
 {}^{\mathrm I}E_1^{p,q}&=H_{\bar\partial}^{p,q}(Y)
          \ \Longrightarrow\ H_{\mathrm{dR}}^{p+q}(Y;\mathbb C),
       &{}^{\mathrm I}d_1&=[\partial],\\
 {}^{\mathrm {II}}E_1^{p,q}&=H_\partial^{p,q}(Y)
          \ \Longrightarrow\ H_{\mathrm{dR}}^{p+q}(Y;\mathbb C),
       &{}^{\mathrm {II}}d_1&=[\bar\partial].
 \end{aligned}
\end{equation}
The second sequence retains the original $(p,q)$ labels, so its
first differential is vertical. Complex conjugation exchanges the
two sequences and reflects the bidegrees. By
Proposition~\ref{hr:prop:frolicher}, both degenerate at their
second page, with only the two corner classes remaining.

The definitions can first be evaluated on elementary complexes.
A \emph{dot} is a one-dimensional space in one bidegree, with
zero differentials. It contributes one class to both quotients.
A \emph{one-arrow complex} has one-dimensional source and target,
joined by either differential as an isomorphism. Its Bott--Chern
class lies at the target, and its Aeppli class lies at the source.
Indeed, the mixed differential is zero. The target is closed for
both differentials, while the Aeppli denominator removes it.

An \emph{elementary square} has basis $a,b,c,e$ in bidegrees
$(p,q),(p+1,q),(p,q+1),(p+1,q+1)$, with
\begin{equation}\label{bc:eq:square}
 \partial a=b,\qquad \bar\partial a=c,\qquad
 \partial c=-e,\qquad \bar\partial b=e.
\end{equation}
All other differentials vanish. The simultaneous kernel is the
line spanned by $e=-\partial\bar\partial a$, so its Bott--Chern
quotient is zero. The mixed kernel is spanned by $b,c,e$.
All three vectors are incoming boundaries, so the Aeppli quotient
also vanishes. This square therefore contributes no class to
either theory.

\subsection{The algebraic double complex and its cohomologies}

A \emph{zigzag} is a double complex whose nonzero spaces are
one-dimensional and whose support is a finite alternating
horizontal and vertical path. Its arrows are isomorphisms, with
successive arrows alternating between a common source and a
common target. Its length is the number of vertices.
Dots and one-arrow complexes have lengths one and two.

Stelzig's structure theorem decomposes every bounded double
complex over a field into elementary squares and zigzags
\cite[Theorem~3]{StelzigDouble}.
The multiplicities are unique, while a splitting generally is not.
The theorem permits arbitrary vector-space dimensions in each
bidegree. Thus it applies to $A$, whose support is bounded.

Let $\mathsf C$ be the direct sum of two dots, at $(0,0)$ and
$(3,3)$, and the following eight one-arrow complexes:
\begin{equation}\label{bc:eq:horizontal-arrows}
 \begin{array}{ll}
 (0,1)\xrightarrow{\ \partial\ }(1,1),&
 (1,1)\xrightarrow{\ \partial\ }(2,1),\\[2pt]
 (1,2)\xrightarrow{\ \partial\ }(2,2),&
 (2,2)\xrightarrow{\ \partial\ }(3,2),
 \end{array}
\end{equation}
and their conjugate reflections
\begin{equation}\label{bc:eq:vertical-arrows}
 \begin{array}{ll}
 (1,0)\xrightarrow{\ \bar\partial\ }(1,1),&
 (1,1)\xrightarrow{\ \bar\partial\ }(1,2),\\[2pt]
 (2,1)\xrightarrow{\ \bar\partial\ }(2,2),&
 (2,2)\xrightarrow{\ \bar\partial\ }(2,3).
 \end{array}
\end{equation}
Each vertex in these displays means a copy of $\mathbb C$, and
each arrow is the identity in the chosen bases. Coincident
bidegrees in different summands contain independent copies.
In particular, consecutive displayed arrows do not compose.
The total vector-space dimension of $\mathsf C$ is $18$.

\begin{theorem}\label{bc:thm:double-complex}
There is an isomorphism of algebraic double complexes
\begin{equation}\label{bc:eq:splitting}
             A\simeq\mathsf C\oplus\mathsf Q,
\end{equation}
where $\mathsf Q$ is a possibly infinite direct sum of elementary
squares. The isomorphism need not be canonical, continuous or
multiplicative.
\end{theorem}

\begin{proof}
An even zigzag of length $2r$ contributes one rank-one
differential on page $r$ of one of the two sequences
\eqref{bc:eq:two-sequences}
\cite[Proposition~6 and Corollary~7]{StelzigDouble}.
Both sequences degenerate at $E_2$. Hence no even zigzag of
length at least four occurs.
For length two, the multiplicity at a specified source is the
rank of the corresponding first differential.
The four ranks in Proposition~\ref{hr:prop:frolicher} give
\eqref{bc:eq:horizontal-arrows}. Conjugation gives precisely
\eqref{bc:eq:vertical-arrows}.

An odd zigzag contributes one total-cohomology class, in the
degree shared by its two endpoints
\cite[Proposition~6]{StelzigDouble}.
The Betti numbers \eqref{hr:eq:betti} restrict this degree to
zero or six. An odd zigzag of length greater than one has
distinct endpoints. Within $0\le p,q\le3$, however, there is
only one position in each of these two total degrees.
Thus the odd zigzags are dots at $(0,0)$ and $(3,3)$.
Their multiplicities are $b_0=b_6=1$.
The structure theorem now leaves only square summands.
\end{proof}

Both sequences are needed in this argument. The four vertical
arrows are acyclic for $\bar\partial$ and for the total
differential, but contribute to Bott--Chern and Aeppli cohomology.
The Betti numbers, together with the support range, exclude
long odd zigzags. Degeneration alone does not exclude them.

\begin{corollary}\label{bc:cor:tables}
With rows indexed by $q$ and columns by $p$, the dimensions are
\begin{equation}\label{bc:eq:tables}
 \begin{array}{c|rrrr}
 h_{BC}^{p,q}&p=0&1&2&3\\ \hline
 q=0&1&0&0&0\\
 1&0&2&1&0\\
 2&0&1&2&1\\
 3&0&0&1&1
 \end{array}
 \qquad
 \begin{array}{c|rrrr}
 h_A^{p,q}&p=0&1&2&3\\ \hline
 q=0&1&1&0&0\\
 1&1&2&1&0\\
 2&0&1&2&0\\
 3&0&0&0&1
 \end{array}.
\end{equation}
\end{corollary}

\begin{proof}
Kernels, images and the displayed quotients commute with direct
sums. The calculation following
\eqref{bc:eq:definition-A} therefore applies to
\eqref{bc:eq:splitting}. Each arrow contributes its target to
Bott--Chern and its source to Aeppli cohomology. Add the two
dots. Every square contributes zero by
\eqref{bc:eq:square}. This gives every entry of
\eqref{bc:eq:tables}, in agreement with the general corner
description in \cite[Lemma~8]{StelzigDouble}.

More explicitly, choose an isomorphism in
\eqref{bc:eq:splitting}. Inclusion of $\mathsf C$ and projection
onto $\mathsf C$ are maps of double complexes. Their composites
are the identity on $\mathsf C$ and the projection with kernel
$\mathsf Q$ on $A$. Since $\mathsf Q$ has zero cohomology in
both theories, these maps induce inverse cohomology isomorphisms.
\end{proof}

The tables satisfy conjugation symmetry and
$h_A^{p,q}=h_{BC}^{3-p,3-q}$. Define the non-$\partial\bar\partial$
degrees by
\[
 \Delta^k(Y)=\sum_{p+q=k}\bigl(h_{BC}^{p,q}+h_A^{p,q}\bigr)-2b_k(Y).
\]
Their values are
\begin{equation}\label{bc:eq:degrees}
       (\Delta^0,\ldots,\Delta^6)=(0,2,4,4,4,2,0).
\end{equation}
The additive comparison does not choose preferred smooth
representatives. In bidegree $(1,1)$, the actual differentials
give a more intrinsic description.

\subsection{Two exact classes in bidegree \texorpdfstring{$(1,1)$}{(1,1)}}

The identity on forms induces projections
\[
 \pi_{\bar\partial}:H_{BC}^{1,1}\longrightarrow H_{\bar\partial}^{1,1},
 \qquad
 \pi_\partial:H_{BC}^{1,1}\longrightarrow H_\partial^{1,1},
\]
and maps from $H_{BC}^{1,1}$ to $H_A^{1,1}$ and
$H_{\mathrm{dR}}^2(Y;\mathbb C)$.
The denominators in \eqref{bc:eq:definition-BC} vanish in each
target, so these maps are well-defined.

\begin{proposition}\label{bc:prop:two-lines}
The differentials induce a canonical isomorphism
\begin{equation}\label{bc:eq:two-lines}
 \begin{aligned}
 H_{\bar\partial}^{0,1}(Y)\oplus H_\partial^{1,0}(Y)
       &\xrightarrow{\ \sim\ }H_{BC}^{1,1}(Y),\\
 ([\gamma],[\eta])&\longmapsto
                       [\partial\gamma+\bar\partial\eta]_{BC}.
 \end{aligned}
\end{equation}
Each summand is one-dimensional.
The kernel of $\pi_{\bar\partial}$ is the second image line,
and the kernel of $\pi_\partial$ is the first.
Both projections have rank one. The maps to Aeppli and
de Rham cohomology are zero.
\end{proposition}

\begin{proof}
For $\bar\partial\gamma=0$, the form $\partial\gamma$ is
closed for both differentials. Replacing $\gamma$ by
$\gamma+\bar\partial f$ changes it by
$\partial\bar\partial f$, which is zero in Bott--Chern
cohomology. The conjugate argument defines the second map.

The composite of the first map with $\pi_{\bar\partial}$ is
the nonzero arrow
${}^{\mathrm I}d_1:H_{\bar\partial}^{0,1}\to
H_{\bar\partial}^{1,1}$. Its source has dimension one.
The conjugate composite with $\pi_\partial$ proves that
the second map also has rank one.
Its image lies in $\ker\pi_{\bar\partial}$, because
$\bar\partial\eta$ is $\bar\partial$-exact.
The two image lines are therefore independent.
Their dimensions sum to $h_{BC}^{1,1}=2$, proving the isomorphism
and both kernel statements.

Each image is an incoming boundary in the Aeppli quotient.
Moreover, $d\gamma=\partial\gamma$ and
$d\eta=\bar\partial\eta$ for the chosen representatives.
Thus both images also vanish in de Rham cohomology.
\end{proof}

In particular, the $\partial\bar\partial$-lemma fails on $Y$.
In its cohomological formulation, this lemma requires the
natural map $H_{BC}^{p,q}\to H_A^{p,q}$ to be injective for
every bidegree. Here that map is zero on the nonzero
two-dimensional space $H_{BC}^{1,1}$.

Wedge product makes $H_{BC}^{\bullet,\bullet}$ a bigraded algebra.
Indeed, wedging a $\partial\bar\partial$-boundary with a form
closed for both differentials gives another
$\partial\bar\partial$-boundary. Its unit is the constant class.
This multiplication comes from actual forms and does not require
a multiplicative splitting of \eqref{bc:eq:splitting}.

\begin{proposition}\label{bc:prop:products}
Choose a nonzero $[\gamma]\in H_{\bar\partial}^{0,1}(Y)$
and a $\bar\partial$-closed representative. Put
\[
             \alpha=\partial\gamma,\qquad
             \beta=\bar\partial\overline\gamma.
\]
Their Bott--Chern classes form a basis of $H_{BC}^{1,1}$ and satisfy
\begin{equation}\label{bc:eq:products}
       [\alpha]_{BC}[\beta]_{BC}=0,\qquad
       [\alpha]_{BC}^3=[\beta]_{BC}^3=0.
\end{equation}
\end{proposition}

\begin{proof}
Conjugation gives $\partial\overline\gamma=0$.
Proposition~\ref{bc:prop:two-lines} supplies the basis.
Both forms are $d$-exact and closed for each differential.
The graded Leibniz rules give
\[
 \bar\partial(\gamma\wedge\overline\gamma)
       =-\gamma\wedge\beta,
 \qquad
 \partial\bar\partial(\gamma\wedge\overline\gamma)
       =-\alpha\wedge\beta.
\]
This proves the mixed-product relation with its sign.

Integration defines an isomorphism
$H_{BC}^{3,3}(Y)\to\mathbb C$.
It is well-defined by Stokes' theorem and nonzero on a positive
top form. Its source has dimension one by
\eqref{bc:eq:tables}. Since
\[
 \alpha^3=d(\gamma\wedge\alpha^2),\qquad
 \beta^3=d(\overline\gamma\wedge\beta^2),
\]
both integrals vanish. The cubic relations follow.
\end{proof}

\subsection{The pullback from the base and the two squares}

The mixed relation follows from the two differentials alone.
To determine the individual squares, we use the holomorphic map
$f:Y\to\mathbb P^1$ of the selected threefold.
Retain the representatives $\gamma,\alpha,\beta$ of
Proposition~\ref{bc:prop:products}. Write
\[
 \mathsf B^{p,q}=H_{BC}^{p,q}(Y),\qquad
 \mathsf B=\bigoplus_{p,q}\mathsf B^{p,q},\qquad
 \mathsf B_+=\bigoplus_{p+q>0}\mathsf B^{p,q}.
\]
These are finite-dimensional complex spaces with the indicated
bigrading. No completion or specialization of a parameter enters
their definition.

\begin{lemma}\label{bc:lem:squares}
The classes $[\alpha]_{BC}$ and $[\beta]_{BC}$ have square zero.
The pullback of a positive Fubini--Study form on $\mathbb P^1$
defines a nonzero real class in $\mathsf B^{1,1}$ with square zero.
\end{lemma}

\begin{proof}
The differentials define two lines
\[
 L=\partial H_{\bar\partial}^{1,2}(Y)\subset\mathsf B^{2,2},
 \qquad
 \overline L=\bar\partial H_\partial^{2,1}(Y)
                                      \subset\mathsf B^{2,2}.
\]
The map defining $L$ is well-defined by the same boundary
calculation as in Proposition~\ref{bc:prop:two-lines}.
Its composite with the Dolbeault projection is the rank-one
first differential from $(1,2)$ to $(2,2)$.
Thus it is injective, since $h^{1,2}=1$.
The conjugate line is nonzero and has zero Dolbeault image.
The two lines are independent, and \eqref{bc:eq:tables} gives
\begin{equation}\label{bc:eq:22-lines}
                    \mathsf B^{2,2}=L\oplus\overline L.
\end{equation}
Moreover,
\[
 \alpha^2=\partial(\gamma\wedge\alpha),\qquad
 \bar\partial(\gamma\wedge\alpha)=0.
\]
Hence $[\alpha^2]_{BC}\in L$, and conjugation gives
$[\beta^2]_{BC}\in\overline L$.

Choose a real positive Fubini--Study form $\omega_{FS}$ with
integral one and put $\eta=f^*\omega_{FS}$.
This is a smooth real semipositive closed $(1,1)$-form.
It is nonzero because $f$ is a submersion on its regular torus
locus. Since the base is a curve,
\begin{equation}\label{bc:eq:base-square}
                         \eta^2=f^*(\omega_{FS}^2)=0.
\end{equation}

Its Bott--Chern class is nonzero. Otherwise
$\eta=\partial\bar\partial h$ for a global smooth complex function.
Reality would give $\eta=i\partial\bar\partial\operatorname{Im}h$.
Thus $\operatorname{Im}h$ would be plurisubharmonic on compact
connected $Y$. The strong maximum principle makes it constant,
contrary to $\eta\ne0$.
One may apply that principle to the trace against any Hermitian
metric, without a K\"ahler hypothesis.

Proposition~\ref{bc:prop:two-lines} and reality now give
\[
 [\eta]_{BC}=c[\alpha]_{BC}+\overline c[\beta]_{BC},
 \qquad c\ne0.
\]
Square this identity. The mixed term vanishes by
\eqref{bc:eq:products}, so \eqref{bc:eq:base-square} yields
\[
      c^2[\alpha^2]_{BC}+\overline c^{\,2}[\beta^2]_{BC}=0.
\]
The two terms belong to the independent lines
\eqref{bc:eq:22-lines}. Each therefore vanishes.
\end{proof}

Rescaling $\gamma$ by $c$ normalizes $[\eta]_{BC}$ to
$[\alpha]_{BC}+[\beta]_{BC}$.
This identifies a geometric semipositive class.
It does not determine its coefficient in a separately prescribed
period normalization or classify all Bott--Chern classes represented
by semipositive closed $(1,1)$-forms.

\subsection{All positive-degree products vanish}

\begin{theorem}\label{bc:thm:ring}
The Bott--Chern algebra of $Y$ is
\begin{equation}\label{bc:eq:ring}
             \mathsf B=\mathbb C\cdot1\oplus\mathsf B_+,
             \qquad \mathsf B_+^2=0.
\end{equation}
Its underlying bigraded vector space is given by
\eqref{bc:eq:tables}.
\end{theorem}

\begin{proof}
Lemma~\ref{bc:lem:squares} and \eqref{bc:eq:products} prove
that every product of two classes in $\mathsf B^{1,1}$ is zero.
Next choose a $\bar\partial$-closed $(1,1)$-form $\zeta$ such
that $v=\partial\zeta$ has nonzero Dolbeault class in $(2,1)$.
The rank-one first differential supplies this choice.
The form $v$ is closed for both differentials and represents
the generator of $\mathsf B^{2,1}$.
Its conjugate generates $\mathsf B^{1,2}$.

The Dolbeault class of $\gamma\wedge\alpha$ maps to
$[\alpha^2]_{BC}=0$ under the injective map defining $L$.
Consequently there is a smooth $(1,1)$-form $\xi$ with
\begin{equation}\label{bc:eq:gamma-alpha}
                   \gamma\wedge\alpha=\bar\partial\xi.
\end{equation}
The natural map
$\mathsf B^{3,2}\to H_{\bar\partial}^{3,2}$ is an isomorphism.
Indeed, both spaces are one-dimensional, and the last rank-one
first differential supplies a nonzero element in its image.
The Dolbeault integration pairing between
$H_{\bar\partial}^{3,2}$ and $H_{\bar\partial}^{0,1}$ is perfect.
The latter space is generated by $[\gamma]$, and
\eqref{bc:eq:gamma-alpha} gives
\[
 \int_Y\gamma\wedge\alpha\wedge v
      =\int_Y\bar\partial(\xi\wedge v)=0.
\]
Thus $[\alpha v]_{BC}=0$.
This use of Dolbeault duality requires only a compact Hermitian
metric: the antilinear Hodge operator gives complementary
harmonic representatives and pairs a nonzero form to its squared
norm. Stokes' theorem makes the pairing independent of representatives.

For the second generator, the mixed primitive is explicit:
\begin{equation}\label{bc:eq:mixed-degree-three}
         \beta\wedge v
             =\partial\bar\partial(\overline\gamma\wedge\zeta).
\end{equation}
Here $\partial\overline\gamma=0$, $\bar\partial\zeta=0$,
and $\partial\beta=0$ give the displayed sign.
Therefore $[\beta v]_{BC}=0$ as well.
Conjugation proves the corresponding products with
$\mathsf B^{1,2}$.

The table \eqref{bc:eq:tables} leaves only five pairs of
positive bidegrees with nonzero possible targets:
\[
 \begin{gathered}
 \mathsf B^{1,1}\mathsf B^{1,1}\subset\mathsf B^{2,2},\qquad
 \mathsf B^{1,1}\mathsf B^{2,1}\subset\mathsf B^{3,2},\qquad
 \mathsf B^{1,1}\mathsf B^{1,2}\subset\mathsf B^{2,3},\\
 \mathsf B^{1,1}\mathsf B^{2,2}\subset\mathsf B^{3,3},\qquad
 \mathsf B^{2,1}\mathsf B^{1,2}\subset\mathsf B^{3,3}.
 \end{gathered}
\]
The first three products have just been computed.
For either remaining pair, both factors are $d$-closed and
have positive total degree below six.
The Betti vanishing \eqref{hr:eq:betti} makes each factor
$d$-exact, so the product has zero integral.
Integration is an isomorphism $\mathsf B^{3,3}\to\mathbb C$,
as proved in Proposition~\ref{bc:prop:products}.
Thus these two products also vanish.
Graded commutativity handles reversed factors.
Every other pair has zero target or leaves $0\le p,q\le3$.
\end{proof}

\subsection{The Aeppli module and integration}

Put $\mathsf A^{p,q}=H_A^{p,q}(Y)$ and
$\mathsf A=\bigoplus_{p,q}\mathsf A^{p,q}$.
Wedge by forms closed for both differentials makes $\mathsf A$
a left $\mathsf B$-module. It does not define an Aeppli ring.
To verify the action, let $a$ satisfy $\partial\bar\partial a=0$.
For a homogeneous form $h$ of total degree $|h|$, the Leibniz
identities give
\begin{equation}\label{bc:eq:Aeppli-boundary}
 (\partial\bar\partial h)\wedge a
 =\partial(\bar\partial h\wedge a)
       +(-1)^{|h|}\bar\partial(h\wedge\partial a).
\end{equation}
Thus changing a Bott--Chern representative changes the product
by an Aeppli boundary. Changing $a$ by a $\partial$- or
$\bar\partial$-boundary also changes the product by such a boundary.
The constant Bott--Chern class acts as the identity.

Integration gives a perfect pairing
\begin{equation}\label{bc:eq:integration-pairing}
 \begin{aligned}
 \mathsf B^{p,q}\times\mathsf A^{3-p,3-q}&\longrightarrow\mathbb C,\\
 (b,a)&\longmapsto\int_Y b\wedge a.
 \end{aligned}
\end{equation}
Stokes' theorem proves that it is well-defined.
For perfection, choose an auxiliary Hermitian metric.
The Bott--Chern and Aeppli Hodge decompositions identify the
two groups with kernels of elliptic operators.
The antilinear Hodge operator exchanges their complementary
kernels and pairs a nonzero representative to its positive squared norm
\cite[Theorem~2.2 and Section~2.c]{SchweitzerBC}.
The resulting pairing is intrinsic.
In particular, there is a unique $\tau_A\in\mathsf A^{3,3}$
with $\int_Y\tau_A=1$.

\begin{proposition}\label{bc:prop:Aeppli-module}
The pairing \eqref{bc:eq:integration-pairing} identifies
\[
                \mathsf A^{p,q}\simeq(\mathsf B^{3-p,3-q})^*.
\]
For homogeneous $x\in\mathsf B_+\cap\mathsf B^{p,q}$ and
$a\in\mathsf A^{r,s}$, the entire module action is
\begin{equation}\label{bc:eq:Aeppli-action}
 x\cdot a=
 \begin{cases}
  \displaystyle\left(\int_Yx\wedge a\right)\tau_A,
                  &(p+r,q+s)=(3,3),\\[4pt]
  0,              &\text{otherwise}.
 \end{cases}
\end{equation}
Together with the identity action of $1\in\mathsf B$, this
determines the module without a choice of bases.
\end{proposition}

\begin{proof}
For homogeneous classes in complementary degrees, integration satisfies
\[
                 \langle b,x\cdot a\rangle=\langle bx,a\rangle.
\]
The factors remain in the displayed order, so no interchange sign
is implicit. If $x\cdot a$ has bidegree other than $(3,3)$,
every class $b$ in its complementary group has positive degree.
Theorem~\ref{bc:thm:ring} gives $bx=0$.
Perfection forces $x\cdot a=0$.
In top bidegree, pairing with the unit gives exactly
\eqref{bc:eq:Aeppli-action}. Outputs outside the support are zero.
\end{proof}

The natural map $H_{BC}^{\bullet,\bullet}\to H_A^{\bullet,\bullet}$
is multiplication by the Aeppli class of $1$.
It is the identity on constants, an isomorphism in bidegree
$(3,3)$ normalized by integration, and zero in every other
positive bidegree.

These results determine multiplication on cohomology and its
integration-dual module. The algebraic splitting
\eqref{bc:eq:splitting} remains nonmultiplicative.
A multiplicative finite model of the smooth double complex,
formality, and higher Bott--Chern--Aeppli operations require
additional constructions. Section~\ref{sec:tangent-curves}
evaluates the Gysin classes of the toric curves by their actual
periods. Their positive exact relations will exclude
pluriclosed Hermitian metrics on the selected threefold.
We first compute the differential class map relevant to these
forms on the same selected threefold.
